\documentclass[
    language=english,
    doctype=syllabus,
    institution=none,
]{../../omnilatex}

\RequirePackage{config/document-settings}

\begin{document}

\maketitle

\section{Course Description}

This course provides a broad introduction to machine learning and statistical
pattern recognition.  Topics include supervised learning (linear regression,
logistic regression, neural networks, support vector machines, decision trees,
and ensemble methods), unsupervised learning (clustering, dimensionality
reduction, and generative models), and the theoretical foundations that underpin
modern machine learning practice.  Students will gain hands-on experience by
implementing algorithms in Python and applying them to real-world datasets.

\section{Learning Objectives}

\begin{objectives}
    \item Formulate supervised and unsupervised learning problems and select
          appropriate models for a given task.
    \item Explain the bias--variance trade-off and apply regularization
          techniques to control model complexity.
    \item Implement and evaluate core algorithms including linear regression,
          logistic regression, decision trees, and $k$-means clustering from
          scratch.
    \item Apply neural network architectures using modern deep learning
          frameworks and diagnose common training issues (vanishing gradients,
          overfitting).
    \item Perform rigorous model evaluation using cross-validation, ROC curves,
          and appropriate performance metrics.
    \item Communicate experimental results clearly in written reports, including
          reproducible code and well-justified conclusions.
\end{objectives}

\section{Required Textbooks and Materials}

\begin{itemize}
    \item Christopher M.~Bishop, \emph{Pattern Recognition and Machine
          Learning}, Springer, 2006. (Primary reference)
    \item Aur\'{e}lien G\'{e}ron, \emph{Hands-On Machine Learning with
          Scikit-Learn, Keras, and TensorFlow}, 3rd~ed., O'Reilly, 2022.
    \item Additional readings and lecture slides will be posted on the course
          website.
\end{itemize}

\noindent\textbf{Software:} Python 3.11+, NumPy, scikit-learn, PyTorch (or
TensorFlow), Jupyter Notebooks.  All required software is available through
the university computing labs or via free open-source installations.

\section{Grading}

Your final grade is determined by the weighted components shown below.
All assignments are graded on a 0--100 scale; letter grades follow the
standard university curve.

\gradingpolicy
\gradingitem{Homework (4 assignments)}{20}
\gradingitem{Programming Projects (2)}{25}
\gradingitem{Midterm Examination}{20}
\gradingitem{Final Examination}{25}
\gradingitem{Participation \& Peer Review}{10}
\endgradingpolicy

\section{Course Policies}

\subsection{Attendance}

Attendance at lectures is expected but not formally recorded.  You are
responsible for all material covered in class, including announcements
and schedule changes.  Attendance at exams is mandatory; makeup exams are
granted only for documented medical or family emergencies.

\subsection{Late Work}

Each student has a total of five \textbf{late days} that may be applied to
any homework or project deadline, at most two per assignment.  Once late days
are exhausted, submissions are penalized 10\% per calendar day late.  No
work is accepted more than seven days after the deadline.

\subsection{Academic Integrity}

All submitted work must represent your own effort unless explicitly authorized
as a collaborative assignment.  You may discuss ideas at a high level but must
write all code and prose independently.  Copying from any source---including
online resources and other students---without proper attribution constitutes
academic dishonesty and will be reported to the Office of Student Conduct.

Violations result in a score of zero on the affected assignment for the first
offense and a failing course grade for subsequent offenses.

\subsection{Accommodations}

Students requiring accommodations should register with the Disability Resource
Center and contact the instructor within the first two weeks of the semester.

\section{Course Schedule}

\small
\begin{tabular}{@{}clp{7.5cm}@{}}
    \toprule
    \textbf{Week} & \textbf{Dates} & \textbf{Topic} \\
    \midrule
    \scheduleentry{1}{Sep 1--3}{Course overview; linear algebra review; probability refresher}
    \scheduleentry{2}{Sep 8--10}{Linear regression; least squares; gradient descent}
    \scheduleentry{3}{Sep 15--17}{Logistic regression; maximum likelihood; regularization}
    \scheduleentry{4}{Sep 22--24}{Bias--variance trade-off; cross-validation; model selection}
    \scheduleentry{5}{Sep 29--Oct 1}{\textit{Midterm 1 review}; neural networks: perceptron \& back-propagation}
    \scheduleentry{6}{Oct 6--8}{Deep learning: activation functions, optimization, CNNs}
    \scheduleentry{7}{Oct 13--15}{Support vector machines; kernel methods}
    \scheduleentry{8}{Oct 20--22}{Decision trees \& random forests; ensemble methods}
    \scheduleentry{9}{Oct 27--29}{Unsupervised learning: $k$-means, hierarchical clustering}
    \scheduleentry{10}{Nov 3--5}{Dimensionality reduction: PCA, t-SNE}
    \scheduleentry{11}{Nov 10--12}{Generative models: GMMs, autoencoders, introduction to GANs}
    \scheduleentry{12}{Nov 17--19}{Reinforcement learning basics; course review; \textit{Final exam review} \\
    \midrule
    & \multicolumn{2}{l}{\textbf{Midterm Examination:} October 8 (in class)} \\
    & \multicolumn{2}{l}{\textbf{Final Examination:} December 10, 9:00--11:00\,AM} \\
    \bottomrule
\end{tabular}

\end{document}
